

%% %%%%%%%%%%%%%%%%%%%%%%%%%%%%%%%%%%%%%%%%%%%%%%%%%
%% Template for a conference paper, prepared for the
%% Food and Resource Economics Department - IFAS
%% UNIVERSITY OF FLORIDA
%% %%%%%%%%%%%%%%%%%%%%%%%%%%%%%%%%%%%%%%%%%%%%%%%%%
%% Version 1.0 // November 2019
%% %%%%%%%%%%%%%%%%%%%%%%%%%%%%%%%%%%%%%%%%%%%%%%%%%
%% Ariel Soto-Caro
%%  - asotocaro@ufl.edu
%%  - arielsotocaro@gmail.com
%% %%%%%%%%%%%%%%%%%%%%%%%%%%%%%%%%%%%%%%%%%%%%%%%%%
\documentclass[12pt]{article}
\usepackage{UF_FRED_paper_style}
\usepackage{threeparttable}
\usepackage{booktabs}
\usepackage{tabularx}
\usepackage{amsmath} 
\usepackage{lipsum}  %% Package to create dummy text (comment or erase before start)
\usepackage{graphicx}

\usepackage{amsthm}
\usepackage[dvipsnames]{xcolor}
\theoremstyle{plain}
\newtheorem{thm}{Theorem}
\newtheorem{defi}{Definition}
\newtheorem{ass}{Assumption}
\newtheorem{prop}{Proposition}
\newtheorem{lem}{Lemma}
\newtheorem{coro}{Corollary}

\usepackage{hyperref}					% Reference
\usepackage[nameinlink]{cleveref} % nameinlink optional but nice
\hypersetup{colorlinks,%				% Reference color setup	
	linkcolor=red,%
	citecolor=blue}
% Tell cleveref how to name this environment
\crefname{thm}{Theorem}{Theorems}
\Crefname{thm}{Theorem}{Theorems}
\crefname{ass}{Assumption}{Assumptions}
\Crefname{ass}{Assumption}{Assumptions}
\crefname{defi}{Definition}{Definitions}
\Crefname{defi}{Definition}{Definitions}
\crefname{prop}{Proposition}{Propositions}
\Crefname{prop}{Proposition}{Propositions}
\crefname{lem}{Lemma}{Lemmas}
\Crefname{lem}{Lemma}{Lemmas}
\crefname{coro}{Corollary}{Corollaries}
\Crefname{coro}{Corollary}{Corollaries}

\newcommand{\red}[1]{{\color{red} #1}}
\newcommand{\blue}[1]{{\color{blue} #1}}
\newcommand{\orange}[1]{{\color{orange} #1}}
\newcommand{\gray}[1]{{\color{gray} #1}}
\newcommand{\green}[1]{{\color{OliveGreen} #1}}
\newcommand{\proptoinverse}{\mathrel{\mskip1mu\reflectbox{$\propto$}\mskip-1mu}}
\usepackage{dcolumn}
\newcolumntype{d}[1]{D..{#1}}
\newcommand{\mc}[1]{\multicolumn{1}{c@{}}{#1}}
\newcommand{\mX}[1]{\multicolumn{1}{X}{\centering #1}}
%% ===============================================
%% Setting the line spacing (3 options: only pick one)
% \doublespacing
% \singlespacing
\onehalfspacing
%% ===============================================

\setlength{\droptitle}{-5em} %% Don't touch

% %%%%%%%%%%%%%%%%%%%%%%%%%%%%%%%%%%%%%%%%%%%%%%%%%%%%%%%%%%
% SET THE TITLE
% %%%%%%%%%%%%%%%%%%%%%%%%%%%%%%%%%%%%%%%%%%%%%%%%%%%%%%%%%%

% TITLE:
\title{Two-country Version of \cite{hirano2025technological}
%\thanks{}
}

% AUTHORS:
\author{
    }
    
% DATE:
\date{\today}

% %%%%%%%%%%%%%%%%%%%%%%%%%%%%%%%%%%%%%%%%%%%%%%%%%%%%%%%%%%
% %%%%%%%%%%%%%%%%%%%%%%%%%%%%%%%%%%%%%%%%%%%%%%%%%%%%%%%%%%
\begin{document}
% %%%%%%%%%%%%%%%%%%%%%%%%%%%%%%%%%%%%%%%%%%%%%%%%%%%%%%%%%%
% %%%%%%%%%%%%%%%%%%%%%%%%%%%%%%%%%%%%%%%%%%%%%%%%%%%%%%%%%%
% ABSTRACT
% %%%%%%%%%%%%%%%%%%%%%%%%%%%%%%%%%%%%%%%%%%%%%%%%%%%%%%%%%%
% %%%%%%%%%%%%%%%%%%%%%%%%%%%%%%%%%%%%%%%%%%%%%%%%%%%%%%%%%%
{\setstretch{.8}
\maketitle
}
\section{Model setup}

\subsection{Representative household}

%--------------------------------------------------------
% US household (equities + US bond only; bond-share cost)
%--------------------------------------------------------
\subsubsection{US}

Let
\[
S_t \equiv Q_{US,t}n_{US,t}+Q_{W,t}n_{W,t},
\qquad
\omega_t \equiv \frac{Q_{US,t}n_{US,t}}{S_t},
\qquad
q_{US,t}\equiv \frac{1}{R_{f,t}}
\]
denote, respectively, the market value of equity purchases, the portfolio weight on US equity within equity
holdings, and the date-$t$ price of the one-period US riskless bond (gross risk-free return $R_{f,t}$).
The US household trades US and RoW equities and the US bond only (it does \emph{not} trade the RoW bond).
It chooses
\[
\{C_{y,t},\,n_{US,t},\,n_{W,t},\,B_{US,t+1}\}
\]
to solve
\[
\max
\ (1-\beta)\log C_{y,t}
+\frac{\beta}{1-\gamma}\log E_t\!\left[(C_{o,t+1})^{1-\gamma}\right]
-\frac{\kappa}{2}(\omega_t-\bar\omega)^2
-\eta\,\Upsilon(\theta_t)
\]
subject to
\begin{align*}
C_{y,t}
&+Q_{US,t}n_{US,t}+Q_{W,t}n_{W,t}
+q_{US,t}B_{US,t+1}
= e_{US,t},\\
C_{o,t+1}
&=
(Q_{US,t+1}+D_{US,t+1})n_{US,t}
+(Q_{W,t+1}+D_{W,t+1})n_{W,t}
+B_{US,t+1}.
\end{align*}
Here $\bar\omega\in[1/2,1]$ is the targeted portfolio weight on US equity (home bias). The term
$\eta\,\Upsilon(\theta_t)$ captures a convex cost in the present-value bond share $\theta_t$,
disciplines large bond positions, and preserves the saving separation. We assume that issuance-cost function $\Upsilon$ satisfies standard regularity and Inada-type conditions.

\begin{ass}[Issuance costs]\label{ass_issuance_costs}
The issuance-cost function $\Upsilon:\,(-\infty,1)\to\mathbb R$ is $C^1$, and satisfies
\[
\lim_{\theta\downarrow -\infty}\Upsilon'(\theta)=-\infty,
\qquad
\Upsilon(\theta)\to+\infty \text{ as } \theta\downarrow -\infty.
\]
\end{ass}

\paragraph{Reformulation in terms of total saving, bond share, and equity share.}
Define the present value of the US bond position and total saving:
\[
b_t\equiv q_{US,t}B_{US,t+1},
\qquad
A_t \equiv S_t+b_t,
\qquad
C_{y,t}=e_{US,t}-A_t.
\]
Define the (present-value) bond share in total saving
\[
\theta_t \equiv \frac{b_t}{A_t},
\qquad
S_t=(1-\theta_t)A_t,
\]
and the gross return on total saving
\[
R_{A,t+1}\equiv (1-\theta_t)R_{p,t+1}+\theta_t R_{f,t},
\qquad
R_{f,t}\equiv \frac{1}{q_{US,t}},
\]
where
\[
R_{p,t+1}=\omega_t R_{US,t+1}+(1-\omega_t)R_{W,t+1},
\qquad
R_{i,t+1}=\frac{Q_{i,t+1}+D_{i,t+1}}{Q_{i,t}},
\ \ i\in\{US,W\}.
\]
Old-age consumption satisfies $C_{o,t+1}=A_tR_{A,t+1}$, so the problem can be written as
\[
\max_{A_t,\omega_t,\theta_t}
\ (1-\beta)\log(e_{US,t}-A_t)
+\frac{\beta}{1-\gamma}\log E_t\!\left[(A_tR_{A,t+1})^{1-\gamma}\right]
-\frac{\kappa}{2}(\omega_t-\bar\omega)^2
-\eta\,\Upsilon(\theta_t).
\]
Under log utility in youth and the risk-sensitive aggregator in old age, the optimal saving choice yields
\begin{align}\label{US_saving_rule_bond_cost}
\frac{1-\beta}{e_{US,t}-A_t}=\frac{\beta}{A_t}
\quad\Rightarrow\quad
A_t=\beta e_{US,t},
\qquad
C_{y,t}=(1-\beta)e_{US,t}.
\end{align}

\paragraph{Portfolio FOCs and stock-pricing implications.}
Define the normalized \emph{undistorted} portfolio kernel
\begin{equation}\label{US_K_def}
\mathcal K_{t,t+1}
\equiv
\frac{R_{A,t+1}^{-\gamma}}{E_t\!\left[R_{A,t+1}^{1-\gamma}\right]},
\qquad
E_t\!\left[\mathcal K_{t,t+1}R_{A,t+1}\right]=1.
\end{equation}
The FOC for $\omega_t$ implies
\begin{align}\label{US_omega_FOC_bond_cost}
\beta(1-\theta_t)\,
E_t\!\left[\mathcal K_{t,t+1}\bigl(R_{US,t+1}-R_{W,t+1}\bigr)\right]
=
\kappa(\omega_t-\bar\omega).
\end{align}
The FOC for $\theta_t$ implies
\begin{align}\label{US_theta_FOC_bond_cost}
\beta\,
E_t\!\left[\mathcal K_{t,t+1}\bigl(R_{f,t}-R_{p,t+1}\bigr)\right]
=
\eta\,\Upsilon'(\theta_t).
\end{align}

For later use, define
\[
\phi_t
\equiv
\frac{\kappa}{\beta(1-\theta_t)}(\omega_t-\bar\omega)
=
\frac{\kappa C_{y,t}}{(1-\beta)S_t}(\omega_t-\bar\omega),
\qquad
\lambda_t
\equiv
\frac{\eta\theta_t}{\beta}\,\Upsilon'(\theta_t).
\]
Using
\[
R_{A,t+1}=(1-\theta_t)R_{p,t+1}+\theta_t R_{f,t}
\]
and \eqref{US_K_def}, the bond-share FOC implies
\begin{align}\label{US_Rp_pricing}
E_t\!\left[\mathcal K_{t,t+1}R_{p,t+1}\right]
=
1-\lambda_t.
\end{align}
Combining \eqref{US_Rp_pricing} with \eqref{US_omega_FOC_bond_cost} yields
\begin{align}
E_t\!\left[\mathcal K_{t,t+1}R_{US,t+1}\right]
&=
1-\lambda_t+\phi_t(1-\omega_t),\label{US_RUS_pricing}\\
E_t\!\left[\mathcal K_{t,t+1}R_{W,t+1}\right]
&=
1-\lambda_t-\phi_t\omega_t.\label{US_RW_pricing}
\end{align}

Therefore the stock-pricing equations can be written as
\begin{align}
Q_{US,t}
&=
E_t\!\left[\widetilde M^{\,US}_{t,t+1}\bigl(Q_{US,t+1}+D_{US,t+1}\bigr)\right],\label{US_AP_bond_cost}\\
Q_{W,t}
&=
E_t\!\left[\widetilde M^{\,W}_{t,t+1}\bigl(Q_{W,t+1}+D_{W,t+1}\bigr)\right],\notag
\end{align}
with
\begin{align*}
\widetilde M^{\,US}_{t,t+1}
&=
\frac{\mathcal K_{t,t+1}}
{1-\lambda_t+\phi_t(1-\omega_t)},\\
\widetilde M^{\,W}_{t,t+1}
&=
\frac{\mathcal K_{t,t+1}}
{1-\lambda_t-\phi_t\omega_t}.
\end{align*}
Thus the kernel that prices the US stock differs from the undistorted kernel $\mathcal K_{t,t+1}$ because of
both the home-bias wedge $\phi_t$ and the bond-share wedge $\lambda_t$.

%--------------------------------------------------------
% RoW household (equities + US bond + RoW bond; convenience yield; NO bond-share costs)
%--------------------------------------------------------
\subsubsection{Rest-of-World (RoW)}

Let
\[
S_t^* \equiv Q_{US,t}n^*_{US,t}+Q_{W,t}n^*_{W,t},
\qquad
\omega_t^* \equiv \frac{Q_{US,t}n^*_{US,t}}{S_t^*},
\qquad
q_{W,t}\equiv \frac{1}{R_{f,t}^W}
\]
where $q_{W,t}$ is the date-$t$ price of the one-period RoW riskless bond (gross return $R_{f,t}^W$).
The RoW household trades both equities, the US bond, and the RoW bond, and chooses
\[
\{C^*_{y,t},\,n^*_{US,t},\,n^*_{W,t},\,B^*_{US,t+1},\,B^*_{W,t+1}\}
\]
to solve
\begin{align*}
\max\quad U_t^*
&=
(1-\beta)\log C^*_{y,t}
+\frac{\beta}{1-\gamma}\log E_t\!\left[\left(C^*_{o,t+1}\right)^{1-\gamma}\right]
-\frac{\kappa}{2}\left(\omega_t^*-\bar\omega^*\right)^2
+\chi\log\!\left(q_{US,t}B^*_{US,t+1}\right),
\end{align*}
subject to
\begin{align*}
C^*_{y,t}
&+Q_{US,t}n^*_{US,t}+Q_{W,t}n^*_{W,t}
+q_{US,t}B^*_{US,t+1}+q_{W,t}B^*_{W,t+1}
= e_{W,t},\\
C^*_{o,t+1}
&=
(Q_{US,t+1}+D_{US,t+1})n^*_{US,t}
+(Q_{W,t+1}+D_{W,t+1})n^*_{W,t}
+B^*_{US,t+1}+B^*_{W,t+1}.
\end{align*}
The term $\chi>0$ generates a convenience yield from holding US bonds and the ``exorbitant privilege'' of US bonds as a consequence.

\paragraph{Stock pricing (RoW effective SDF).}
The FOCs for equity positions imply, for $i\in\{US,W\}$,
\[
Q_{i,t}
=
E_t\!\left[\tilde M^{*,i}_{t,t+1}\bigl(Q_{i,t+1}+D_{i,t+1}\bigr)\right],
\]
with
\begin{align*}
\tilde M^{*,US}_{t,t+1}=\frac{M^*_{t,t+1}}{1+\phi_t^*(1-\omega_t^*)},
\qquad
\tilde M^{*,W}_{t,t+1}=\frac{M^*_{t,t+1}}{1-\phi_t^*\omega_t^*},\\
M^*_{t,t+1}
=\frac{\beta}{1-\beta}\,
\frac{C^*_{y,t}(C^*_{o,t+1})^{-\gamma}}{E_t\!\left[(C^*_{o,t+1})^{1-\gamma}\right]},
\qquad
\phi_t^* \equiv \frac{\kappa C^*_{y,t}}{(1-\beta)S_t^*}\,(\omega_t^*-\bar\omega^*).
\end{align*}

\paragraph{Saving notation and portfolio shares.}
Define present values and total saving:
\[
b_{US,t}^*\equiv q_{US,t}B^*_{US,t+1},
\qquad
b_{W,t}^*\equiv q_{W,t}B^*_{W,t+1},
\qquad
A_t^*\equiv S_t^*+b_{US,t}^*+b_{W,t}^*,
\qquad
C^*_{y,t}=e_{W,t}-A_t^*.
\]
Define bond shares in total saving:
\[
\theta_{US,t}^*\equiv \frac{b_{US,t}^*}{A_t^*},
\qquad
\theta_{W,t}^*\equiv \frac{b_{W,t}^*}{A_t^*},
\qquad
S_t^*=(1-\theta_{US,t}^*-\theta_{W,t}^*)A_t^*.
\]
Define the gross return on total saving:
\[
R_{A,t+1}^*
\equiv
(1-\theta_{US,t}^*-\theta_{W,t}^*)R_{p,t+1}^*
+\theta_{US,t}^* R_{f,t}
+\theta_{W,t}^* R_{f,t}^W,
\]
where
\[
R_{f,t}\equiv \frac{1}{q_{US,t}},
\qquad
R_{f,t}^W\equiv \frac{1}{q_{W,t}},
\qquad
R_{p,t+1}^*=\omega_t^*R_{US,t+1}+(1-\omega_t^*)R_{W,t+1}.
\]
Then $C^*_{o,t+1}=A_t^*R_{A,t+1}^*$.

\paragraph{RoW saving rule.}
Because $\chi\log(q_{US,t}B_{US,t+1}^*)=\chi\log(\theta_{US,t}^*A_t^*)$ contributes $\chi\log A_t^*$,
the saving choice separates and yields
\begin{align}\label{RoW_saving_rule_bond_cost}
\frac{1-\beta}{e_{W,t}-A_t^*}=\frac{\beta+\chi}{A_t^*}
\quad\Rightarrow\quad
A_t^*=\frac{\beta+\chi}{1+\chi}\,e_{W,t},
\qquad
C^*_{y,t}=\frac{1-\beta}{1+\chi}\,e_{W,t}.
\end{align}

\paragraph{RoW portfolio FOCs.}
Define the normalized RoW portfolio kernel
\[
\mathcal M_{t,t+1}^* \equiv \frac{(R_{A,t+1}^*)^{-\gamma}}{E_t\!\left[(R_{A,t+1}^*)^{1-\gamma}\right]}.
\]
The FOC for $\omega_t^*$ is
\begin{align}\label{RoW_omega_FOC_mod}
\beta\bigl(1-\theta_{US,t}^*-\theta_{W,t}^*\bigr)\,
E_t\!\left[\mathcal M_{t,t+1}^*\bigl(R_{US,t+1}-R_{W,t+1}\bigr)\right]
=
\kappa(\omega_t^*-\bar\omega^*).
\end{align}
The FOC for the RoW bond share $\theta_{W,t}^*$ is
\begin{align}\label{RoW_thetaW_FOC_mod}
\beta\,E_t\!\left[\mathcal M_{t,t+1}^*\bigl(R_{f,t}^W-R_{p,t+1}^*\bigr)\right]
=
0.
\end{align}
The FOC for the US bond share $\theta_{US,t}^*$ is
\begin{align}\label{RoW_thetaUS_FOC_mod}
\beta\,E_t\!\left[\mathcal M_{t,t+1}^*\bigl(R_{f,t}-R_{p,t+1}^*\bigr)\right]
+\frac{\chi}{\theta_{US,t}^*}
=
0.
\end{align}
And the FOCs for the bond amounts $(B_{W,t+1}^*,B_{US,t+1}^*)$ imply the bond-pricing equations for RoW agents:
\begin{align}
q_{W,t}&=E_t\!\left[M^*_{t,t+1}\right], \label{RoW_BW_FOC}\\
q_{US,t}&=E_t\!\left[M^*_{t,t+1}\right]
+\frac{\chi}{1+\chi}\,\frac{e_{W,t}}{B^*_{US,t+1}}. \label{RoW_BUS_FOC}
\end{align}
Since the convenience-yield term $\chi\log(q_{US,t}B_{US,t+1}^*)$ requires $B_{US,t+1}^*>0$, the wedge term in
\eqref{RoW_BUS_FOC} is strictly positive; combing with \eqref{RoW_BW_FOC} immediately gives  $q_{US,t}>q_{W,t}$, generating the exorbitant privilege of US bonds.



%--------------------------------------------------------
% Market clearing (two bonds with zero net supply)
%--------------------------------------------------------
\subsection{Market clearing}

\paragraph{World stock market clearing.}
\begin{align}\label{stock_clear_mod}
\begin{split}
n_{US,t}+n^*_{US,t}=1,\\
n_{W,t}+n^*_{W,t}=1.
\end{split}
\end{align}

\paragraph{Bond market clearing.}
Assume both one-period riskless bonds are in zero net supply. Since only RoW trades the RoW bond, bond market
clearing is
\begin{align}\label{bond_clear_mod_zerosupply}
B_{US,t+1}+B^*_{US,t+1}=0,
\qquad
B^*_{W,t+1}=0.
\end{align}
Equivalently, in present values,
\[
b_t+b_{US,t}^*=0,
\qquad
b_{W,t}^*=0.
\]

\paragraph{Stock market clearing in value form.}
Using portfolio weights,
\begin{align}\label{stock_weight_clear_mod}
\begin{split}
Q_{US,t} &= \omega_t S_t + \omega_t^* S_t^*,\\
Q_{W,t}  &= (1-\omega_t) S_t + (1-\omega_t^*) S_t^* .
\end{split}
\end{align}

\paragraph{Aggregate market-cap identity.}
Since $S_t=A_t-b_t$ and $S_t^*=A_t^*-b_{US,t}^*-b_{W,t}^*$, we have
\[
Q_{US,t}+Q_{W,t}=S_t+S_t^*=A_t+A_t^*-(b_t+b_{US,t}^*)-b_{W,t}^*.
\]
Using bond clearing \eqref{bond_clear_mod_zerosupply} and the saving rules
\eqref{US_saving_rule_bond_cost} and \eqref{RoW_saving_rule_bond_cost}, it follows that
\begin{align}\label{Qsum_identity_mod_zerosupply}
\boxed{
Q_{US,t}+Q_{W,t}
=
\beta e_{US,t}
+\frac{\beta+\chi}{1+\chi}\,e_{W,t}.
}
\end{align}

\paragraph{Goods market clearing.}
\begin{align}\label{goods_clear_mod}
C_{y,t}+C^*_{y,t}+C_{o,t}+C^*_{o,t}
=
e_{US,t}+e_{W,t}+D_{US,t}+D_{W,t}.
\end{align}

\section{Existence of bubbles}\label{sec:existence_of_bubbles}

In this section we reserve $\mathcal K_{t,t+1}$ for the undistorted normalized portfolio kernel,
\[
\mathcal K_{t,t+1}
=
\frac{R_{A,t+1}^{-\gamma}}{E_t\!\left[R_{A,t+1}^{1-\gamma}\right]},
\]
and we let $\mathcal M_{t,t+1}$ denote the \emph{effective} kernel that prices the US stock. By
\eqref{US_AP_bond_cost},
\begin{equation}\label{effective_US_kernel}
\mathcal M_{t,t+1}
\equiv
\widetilde M^{\,US}_{t,t+1}
=
\frac{\mathcal K_{t,t+1}}{\Psi_t},
\qquad
\Psi_t
\equiv
1-\lambda_t+\phi_t(1-\omega_t),
\end{equation}
where
\[
\phi_t=\frac{\kappa}{\beta(1-\theta_t)}(\omega_t-\bar\omega),
\qquad
\lambda_t=\frac{\eta\theta_t}{\beta}\,\Upsilon'(\theta_t).
\]

By \eqref{US_saving_rule_bond_cost} and the identity $C_{o,t+1}=A_tR_{A,t+1}$, the undistorted
intertemporal marginal rate of substitution equals
\[
\frac{\beta}{1-\beta}\,
\frac{C_{y,t}C_{o,t+1}^{-\gamma}}{E_t\!\left[C_{o,t+1}^{1-\gamma}\right]}
=
\mathcal K_{t,t+1}.
\]
Hence the only difference between the stock-pricing kernel $\mathcal M_{t,t+1}$ and the
undistorted kernel $\mathcal K_{t,t+1}$ is the time-$t$ portfolio wedge $\Psi_t^{-1}$.

We impose the following regularity condition.

\begin{ass}[Regular effective kernel]\label{ass_regular_kernel}
\[
\Psi_t>0
\qquad\text{a.s. for all }t.
\]
\end{ass}

Also note that because the RoW utility includes $\chi\log(q_{US,t}B^*_{US,t+1})$, feasibility requires
$B^*_{US,t+1}>0$. With zero net supply of US bonds,
$B_{US,t+1}+B^*_{US,t+1}=0$ implies $B_{US,t+1}<0$, hence
\begin{equation}\label{theta_nonpos}
\theta_t=\frac{b_t}{A_t}<0
\quad\Rightarrow\quad
1-\theta_t>1.
\end{equation}

Then we show that $\Psi_t>0$ are very likely to hold under mild parameter restrictions.
Because $0\le \omega_t\le 1$, we have $1-\omega_t\in[0,1]$ and
\[
(\omega_t-\bar\omega)(1-\omega_t)
\ge -\bar\omega(1-\omega_t)
\ge -\bar\omega.
\]
Therefore
\[
\phi_t(1-\omega_t)
=
\frac{\kappa}{\beta(1-\theta_t)}(\omega_t-\bar\omega)(1-\omega_t)
\ge
-\frac{\kappa\bar\omega}{\beta(1-\theta_t)}.
\]
We further assume the function form of issuance costs of bonds, for now, as
 $\Upsilon(\theta_t)=-\log(1-\theta_t)+\theta_t^2/2$. Then $\Psi_t >0$ if
\begin{equation} \label{suff_condition_pos_wedge}
	\beta > \frac{\eta \theta_t}{1-\theta_t}+\eta \theta_t^2 +\frac{\kappa\bar\omega}{1-\theta_t},	
\end{equation}
as
\[
\Psi_t > 1-\frac{\eta \theta_t}{\beta} \Upsilon'(\theta_t)-\frac{\kappa\bar\omega}{\beta(1-\theta_t)}.
\]
It is immediate that \eqref{suff_condition_pos_wedge} easily hold for sufficiently high $\beta$ and 
sufficiently low $\eta$ and $\kappa$ as $\theta_t$ usually falls within $[-1,0)$ and $1-\theta_t>1$.

\subsection{Temporary unbalanced growth and stochastic bubbles}

Following \citet{hirano2025technological}, let
\[
\mathcal{M}_{t,t+s}:=\prod_{j=0}^{s-1} \mathcal{M}_{t+j,t+j+1}.
\]
The US stock pricing equation \eqref{US_AP_bond_cost} (for $i=US$) implies
\[
Q_{US,0}=
\mathrm{E}_0 \left[\sum_{s=1}^{\infty} \mathcal{M}_{0,s} D_{US,s}\right]
+\lim_{t \rightarrow \infty} \mathrm{E}_0\left[\mathcal{M}_{0,t} Q_{US,t}\right],
\]
or more generally,
\[
Q_{US,t}=
\underbrace{\mathrm{E}_t \left[ \sum_{s=1}^{\infty} \mathcal{M}_{t,t+s} D_{US,t+s}\right]}_{\text{fundamental value }V_{US,t}}
+\underbrace{\lim_{T \rightarrow \infty} \mathrm{E}_t\left[\mathcal{M}_{t,T} Q_{US,T}\right]}_{\text{bubble }B_{US,t}}.
\]

Consider a regime-switching model with an absorbing state, as in \citet{weil1987confidence}.

\begin{ass}\label{ass_regime_switch}
There are two states of the US economy denoted by $u, b$. Letting $z_t \in \{u, b\}$ denote the state at time $t$,
the transition probabilities are given by
\[
P[z_{t+1} = u \mid z_t = u] = \pi \in (0,1), \qquad P[z_{t+1} = b \mid z_t = b] = 1.
\]
\end{ass}

\Cref{ass_regime_switch} implies that state $b$ is absorbing. We further assume that once state $b$ is reached, uncertainty is resolved,
and endowments and dividends grow at the same rate (balanced growth). On the other hand, the RoW is always on its balanced growth path (BGP hereafter).

\begin{ass}\label{ass_balanced_growth}
For any $\tau$, conditional on $z_\tau = b$, the sequence
$\{e_{US,t}, D_{US,t}, e_{W,t}, D_{W,t}\}_{t=\tau}^\infty$ is deterministic and
\[
\frac{e_{US,t+1}}{e_{US, t}} =
\frac{D_{US, t+1}}{D_{US,t}}=
\frac{e_{W,t+1}}{e_{W, t}} =
\frac{D_{W, t+1}}{D_{W,t}} ,\quad \forall t \geq \tau.
\]
The ratios are well defined, and \(D_{US,\tau}/e_{US,\tau}>0\).
\end{ass}

We can show that once state $b$ is reached, there is no bubble in the US stock market.

\begin{prop}\label{prop1}
If \Cref{ass_regular_kernel,ass_regime_switch,ass_balanced_growth} hold, once state $b$ is reached, the US stock price equals its fundamental value:
$Q_{US,t}=V_{US,t}$.
\end{prop}

\begin{proof}
Fix a date $\tau$ with $z_\tau=b$. By \Cref{ass_regime_switch}, the state remains
$b$ forever, so the US pricing equation is deterministic:
\[
Q_{US,s}
=
\mathcal M_{s,s+1}\bigl(Q_{US,s+1}+D_{US,s+1}\bigr),
\qquad s\ge \tau.
\]
Hence
\[
\mathcal M_{\tau,t}Q_{US,t}
=
Q_{US,\tau}
\prod_{s=\tau+1}^{t}
\frac{Q_{US,s}}{Q_{US,s}+D_{US,s}}.
\]
Balanced growth implies that
\[
\frac{e_{W,s}}{e_{US,s}}=c_W^b<\infty,
\qquad
\frac{D_{US,s}}{e_{US,s}}=c_D^b>0,
\qquad s\ge \tau.
\]
Using \eqref{Qsum_identity_mod_zerosupply},
\[
Q_{US,s}
\le
Q_{US,s}+Q_{W,s}
=
\left(\beta+\frac{\beta+\chi}{1+\chi}c_W^b\right)e_{US,s}
\equiv C_Q^b e_{US,s}.
\]
Therefore
\[
\frac{D_{US,s}}{Q_{US,s}}
\ge
\frac{c_D^b}{C_Q^b}
\equiv \underline d^b>0.
\]
It follows that
\[
0\le
\mathcal M_{\tau,t}Q_{US,t}
\le
Q_{US,\tau}
\left(\frac{1}{1+\underline d^b}\right)^{t-\tau}
\to0.
\]
Thus the transversality term in state $b$ vanishes, so the US stock price
equals its fundamental value.
\end{proof}

By \Cref{prop1}, bubbles cannot arise in state $b$. So we assume $z_0=u$ from now on. Let $\tau$ be the stopping time that represents the first time the US economy switches to state $b$:
$\tau =\min \{t:z_t=b\}$. By \Cref{ass_regime_switch}, $\tau$ follows a geometric distribution,
\[
\mathrm{P}[\tau=j]=\pi^{j-1}(1-\pi), \quad j=1,2,\dots.
\]
Then the bubble term can be written as
\[
\mathrm{E}_0\left[\mathcal{M}_{0 , t} Q_{US,t}\right]
=\sum_{j=1}^t \pi^{j-1}(1-\pi) \mathrm{E}_0\left[\mathcal{M}_{0,t} Q_{US,t} \mid \tau=j\right]
+\pi^t \mathrm{E}_0\left[\mathcal{M}_{0, t} Q_{US,t} \mid \tau>t\right].
\]

The following proposition provides a sufficient condition for the nonexistence of bubbles.

\begin{prop}\label{prop2}
Suppose \Cref{ass_regular_kernel,ass_regime_switch,ass_balanced_growth} hold. Then there is no bubble if and only if
\[
\lim _{t \rightarrow \infty} \pi^t \mathrm{E}_0\left[\mathcal{M}_{0,t} Q_{US,t} \mid \tau>t\right]=0.
\]
\end{prop}

\begin{proof}
Decompose the bubble term according to the first switching date:
\[
\mathrm{E}_0\left[\mathcal{M}_{0 , t} Q_{US,t}\right]
=
\sum_{j=1}^t \pi^{j-1}(1-\pi)
\mathrm{E}_0\left[\mathcal{M}_{0,t} Q_{US,t} \mid \tau=j\right]
+\pi^t \mathrm{E}_0\left[\mathcal{M}_{0, t} Q_{US,t} \mid \tau>t\right].
\]
Conditional on $\tau=j$, the economy is in deterministic balanced growth from
date $j$ onward. By \Cref{prop1},
\[
\lim_{t\to\infty}
\mathrm{E}_0\left[\mathcal M_{0,t}Q_{US,t}\mid \tau=j\right]=0.
\]
Equivalently, on $\{\tau\le t\}$,
\[
\mathcal M_{0,t}Q_{US,t}
=
\mathcal M_{0,\tau}\mathcal M_{\tau,t}Q_{US,t}
\to0.
\]
The geometric bound in the proof of \Cref{prop1} gives
\[
0\le
\mathcal M_{0,t}Q_{US,t}\mathbf 1_{\{\tau\le t\}}
\le
\mathcal M_{0,\tau}Q_{US,\tau},
\]
and the right-hand side is integrable by the stock-pricing equation and finite
date prices. Hence the post-switch component converges to zero by dominated
convergence. Therefore the only possible surviving term is the no-switch
history:
\[
\lim_{t\to\infty}\mathrm{E}_0\left[\mathcal M_{0,t}Q_{US,t}\right]=0
\quad\Longleftrightarrow\quad
\lim_{t\to\infty}
\pi^t \mathrm{E}_0\left[\mathcal M_{0,t}Q_{US,t}\mid \tau>t\right]=0.
\]
\end{proof}

Aligned with the sufficient conditions in \Cref{prop2}, we introduce the following assumption.

\begin{ass}\label{ass_switch_dependence}
Conditional on time $t-1$ information, the US endowment $e_{US,t}$ and US dividend $D_{US,t}$ depend only on the state
$z_t \in \{u,b\}$.
\end{ass}


Now we are ready to present a more general but less primitive version of Theorem 1 in \citet{hirano2025technological}.
Given the regime history, all exogenous fundamentals are deterministic, and we
focus on the corresponding Markov equilibrium in which prices, portfolios, and
consumptions are functions of the regime history.
For $z\in\{u,b\}$, let $x_t^z$ be the value of $x_{US,t}^z$ conditional on $z_0=\cdots=z_{t-1}=u$ and $z_t=z$,
and let $C_t^z=C_{o,t}^z$ denote US old-age consumption at date $t$ on branch $z$.

\begin{thm}\label{thm_1}
Suppose \Cref{ass_regular_kernel,ass_regime_switch,ass_balanced_growth,ass_switch_dependence} hold. If $z_0=u$, then there is a bubble at $t=0$ if and only if
\begin{subequations}\label{thm_cond1}
\begin{align}
\sum_{t=1}^\infty \frac{D_t^u}{Q_t^u}&<\infty, \label{thm_cond_1a}\\
\sum_{t=1}^\infty
{\frac{Q_t^b+D_t^b}{(C_t^b)^\gamma}}\big/\left({\frac{Q_t^u}{(C_t^u)^\gamma}}\right)&<\infty. \label{thm_cond_1b}
\end{align}
\end{subequations}
\end{thm}
\begin{proof}
Let
\[
\mathcal B_t \equiv \mathbb E_0\!\left[\mathcal M_{0,t}\,Q_t^u\,\mathbf 1\{\tau>t\}\right].
\]
Because $\mathbb P(\tau>t)=\pi^t$, this is equivalent to
\[
\mathcal B_t
=
\pi^t\,\mathbb E_0\!\left[\mathcal M_{0,t}Q_t^u\mid \tau>t\right].
\]
By \Cref{prop2}, there is a bubble at $t=0$ if and only if
\[
\lim_{t\to\infty}\mathcal B_t>0,
\]
and there is no bubble if and only if $\lim_{t\to\infty}\mathcal B_t=0$.

On histories with $z_{t-1}=u$, the one-step US stock pricing equation is
\[
Q_{t-1}^u
=
\mathbb E_{t-1}\!\left[
\mathcal M_{t-1,t}
\Bigl(
(Q_t^u+D_t^u)\mathbf 1_{\{z_t=u\}}
+
(Q_t^b+D_t^b)\mathbf 1_{\{z_t=b\}}
\Bigr)
\right].
\]
Multiply both sides by $\mathcal M_{0,t-1}\mathbf 1\{\tau>t-1\}$ and take $\mathbb E_0$. Using
$\mathcal M_{0,t}=\mathcal M_{0,t-1}\mathcal M_{t-1,t}$, we obtain
\begin{equation}\label{recursion_bubble_1}
\mathbb E_0\!\left[\mathcal M_{0,t-1}Q_{t-1}^u\mathbf 1\{\tau>t-1\}\right]
=
\mathbb E_0\!\left[\mathcal M_{0,t}(Q_t^u+D_t^u)\mathbf 1\{\tau>t\}\right]
+
\mathbb E_0\!\left[\mathcal M_{0,t}(Q_t^b+D_t^b)\mathbf 1\{\tau=t\}\right].
\end{equation}
Since
\[
\mathbb E_0\!\left[\mathcal M_{0,t}(Q_t^u+D_t^u)\mathbf 1\{\tau>t\}\right]
=
\mathcal B_t+\mathbb E_0\!\left[\mathcal M_{0,t}D_t^u\mathbf 1\{\tau>t\}\right],
\]
\eqref{recursion_bubble_1} becomes
\begin{equation}\label{bubble_recursion_2}
\mathcal B_{t-1}
=
\mathcal B_t
+\underbrace{\mathbb E_0\!\left[\mathcal M_{0,t}D_t^u\mathbf 1\{\tau>t\}\right]}_{\text{dividend leakage term}}
+\underbrace{\mathbb E_0\!\left[\mathcal M_{0,t}(Q_t^b+D_t^b)\mathbf 1\{\tau=t\}\right]}_{
	\text{switch-at-t leakage term}},
\end{equation}
where the last two terms can be interpreted as leakage terms that prevent the bubble from growing too fast. 
In other words, \textit{leakage} is defined as the amount of present value that drains out of the 
continuation bubble mass when we move one period forward. It’s exactly the gap $\mathcal  B_{t-1}-\mathcal B_t$.

Define the leakage ratio
\[
a_t\equiv
\frac{
\mathbb E_0\!\left[\mathcal M_{0,t}D_t^u\mathbf 1\{\tau>t\}\right]
+
\mathbb E_0\!\left[\mathcal M_{0,t}(Q_t^b+D_t^b)\mathbf 1\{\tau=t\}\right]
}{
\mathcal B_t
}.
\]
Then \eqref{bubble_recursion_2} implies
\[
\mathcal B_{t-1}=(1+a_t)\mathcal B_t,
\]
so iterating yields
\[
\mathcal B_t=\mathcal B_0\prod_{s=1}^t\frac{1}{1+a_s},
\qquad
\mathcal B_0=Q_0^u>0.
\]
Since $a_s\ge 0$,
\[
\mathcal B_0\left(1+\sum_{s=1}^t a_s\right)^{-1}
\ge
\mathcal B_t
\ge
\mathcal B_0\exp\!\left(-\sum_{s=1}^t a_s\right).
\]
Hence
\[
\lim_{t\to\infty}\mathcal B_t>0
\quad\Longleftrightarrow\quad
\sum_{t=1}^\infty a_t<\infty.
\]

It remains to express $a_t$ in terms of branch variables. On $\{\tau>t\}$ we have
$z_1=\cdots=z_t=u$, so
\[
\frac{
\mathbb E_0\!\left[\mathcal M_{0,t}D_t^u\mathbf 1\{\tau>t\}\right]
}{
\mathcal B_t
}
=
\frac{D_t^u}{Q_t^u}.
\]

On $\{\tau=t\}$ we have $z_1=\cdots=z_{t-1}=u$ and $z_t=b$, hence
\[
\mathbb E_0\!\left[\mathcal M_{0,t}(Q_t^b+D_t^b)\mathbf 1\{\tau=t\}\right]
=
\pi^{t-1}(1-\pi)\,
\mathcal M_{0,t-1}^u\,
\mathcal M_{t-1,t}^b\,
(Q_t^b+D_t^b),
\]
while
\[
\mathcal B_t
=
\pi^t\,
\mathcal M_{0,t-1}^u\,
\mathcal M_{t-1,t}^u\,
Q_t^u.
\]
Therefore,
\[
\frac{
\mathbb E_0\!\left[\mathcal M_{0,t}(Q_t^b+D_t^b)\mathbf 1\{\tau=t\}\right]
}{
\mathcal B_t
}
=
\frac{1-\pi}{\pi}\,
\frac{\mathcal M_{t-1,t}^b}{\mathcal M_{t-1,t}^u}\,
\frac{Q_t^b+D_t^b}{Q_t^u},
\]
which is the ratio of the discounted b-state cum-dividend stock price to the discounted u-state ex-dividend stock price.

Finally, under \Cref{ass_switch_dependence}, the time-$(t-1)$ portfolio choices
$(\omega_{t-1},\theta_{t-1})$ are $\mathcal F_{t-1}$-measurable and do not depend on $z_t$.
Hence the time-$(t-1)$ wedge
\[
\Psi_{t-1}=1-\lambda_{t-1}+\phi_{t-1}(1-\omega_{t-1})
\]
is common across the two realizations $z_t\in\{u,b\}$, so it cancels in the ratio:
\[
\frac{\mathcal M_{t-1,t}^b}{\mathcal M_{t-1,t}^u}
=
\frac{\mathcal K_{t-1,t}^b}{\mathcal K_{t-1,t}^u}.
\]
Moreover, the denominator
$E_{t-1}[R_{A,t}^{1-\gamma}]$ in $\mathcal K_{t-1,t}$ is also common across the two realizations,
so
\[
\frac{\mathcal K_{t-1,t}^b}{\mathcal K_{t-1,t}^u}
=
\left(\frac{R_{A,t}^b}{R_{A,t}^u}\right)^{-\gamma}
=
\left(\frac{C_t^b}{C_t^u}\right)^{-\gamma}
\]
Therefore,
\[
a_t
=
\frac{D_t^u}{Q_t^u}
+\frac{1-\pi}{\pi}
\left(\frac{C_t^u}{C_t^b}\right)^\gamma
\frac{Q_t^b+D_t^b}{Q_t^u}
=
\frac{D_t^u}{Q_t^u}
+\frac{1-\pi}{\pi}
\frac{\frac{Q_t^b+D_t^b}{(C_t^b)^\gamma}}{\frac{Q_t^u}{(C_t^u)^\gamma}}.
\]
Hence $\sum_{t=1}^\infty a_t<\infty$ is equivalent to \eqref{thm_cond1}.
\end{proof}

Intuitively, the first summability condition \eqref{thm_cond_1a} requires that dividends along the $u$-history 
do not grow too fast relative to stock prices, while the second summability condition \eqref{thm_cond_1b} requires that the stock prices discounted by the consumption-based SDF along the $u$-history grow much faster than the cum-dividend stock prices would have in b-state. Therefore, the bubble candidate survives if and only if,
 over time, dividends and switching events do not drain too much value relative to what remains; that is, 
 if and only if the cumulative leakage $\sum a_t$ is finite.

\subsection{Primitive sufficient and necessary conditions}

We now derive primitive conditions for the existence of bubbles by translating the two summability conditions
in \Cref{thm_1} into exogenous variables along the $u$-history and at the switching date. We first impose
primitive boundedness restrictions on relative country size, switch-date US fundamentals, portfolio weights,
US dividends, and continuation-branch growth.

\begin{ass}[Relative country size and dividend bounds]\label{ass_relative_dividend_bounds}
There exists $\overline\lambda_{e_W}<\infty$ and $\overline\lambda_{D_W}<\infty$ such that
\[
\frac{e_{W,t}}{e_{US,t}^z}\le \overline\lambda_{e_W},
\qquad
\frac{D_{W,t}}{e_{US,t}^z}\le \overline\lambda_{D_W},
\qquad\text{for all }t\le \tau.
\]
\end{ass}

\begin{ass}[Switch-date US fundamentals]\label{ass_switch_us_fundamentals}
There exist constants $0< \overline\lambda_e<\infty$ and
$0\le \overline\lambda_{D^b_{US}}<\infty$ such that, whenever $z_{t-1}=u$ and $z_t=b$,
\[
\frac{e_{US,t}^b}{e_{US,t}^u}
\le
\overline\lambda_e,
\qquad
\frac{D_{US,t}^b}{e_{US,t}^u}
\le
\overline\lambda_{D^b_{US}}.
\]
\end{ass}

\begin{ass}[Uniformly interior equity weights]\label{ass_interior_equity_weights}
There exists $\underline\lambda_\omega\in(0,1/2)$ such that, for each
$z\in\{u,b\}$ and all $t\le \tau$,
\[
\underline\lambda_\omega
\le
\omega_t^z,\ \omega_t^{*,z}
\le
1-\underline\lambda_\omega.
\]
\end{ass}

\begin{ass}[US dividends and continuation-branch US growth]\label{ass_u_dividend_growth}
There exist constants $0<\underline g_u\le \overline g_u<\infty$ and
$\overline\lambda_{D_{US}}<\infty$ such that
\[
\underline g_u
\le
\frac{e_{US,t}^u}{e_{US,t-1}^u}
\le
\overline g_u,
\qquad
\frac{D_{US,t}^z}{e_{US,t}^z}
\le
\overline\lambda_{D_{US}},
\qquad z\in\{u,b\},\ t\ge 1.
\]
\end{ass}

\begin{ass}[Risk aversion below unity]\label{ass_gamma_lt_one}
\[
0<\gamma<1.
\]
\end{ass}

Then in the following lemmas, we show that stock prices are bounded and proportional to endowments; thus, the return rates for stocks, bonds, and the mixed portfolio are also finite, which finally leads to consumption that is proportional to endowments.

\begin{lem}[Equilibrium bounds on stock market values]\label{lem_stock_value_bounds}
Suppose
\Cref{ass_relative_dividend_bounds,ass_interior_equity_weights}
hold. Then there exist constants $0<\underline q\le \overline q<\infty$ such that, for each
$i\in\{US,W\}$, each $z\in\{u,b\}$, and all $t\le\tau$,
\[
\underline q\,e_{US,t}^z
\le
Q_{i,t}^z
\le
\overline q\,e_{US,t}^z,
\]
where one may take
\[
\underline q\equiv \underline\lambda_\omega\,\beta,
\qquad
\overline q\equiv
\beta
+\frac{\beta+\chi}{1+\chi}\overline\lambda_{e_W}.
\]
\end{lem}

\begin{proof}
For $z=u$, \eqref{US_saving_rule_bond_cost}, \eqref{stock_weight_clear_mod},  \eqref{theta_nonpos}, and
\Cref{ass_interior_equity_weights} imply
\[
Q_{US,t}^u
=
\omega_t^u S_t^u+\omega_t^{*,u}S_t^{*,u}
\ge
\omega_t^u S_t^u
\ge
\underline\lambda_\omega(1-\theta_t^u)A_t^u
\ge
\underline\lambda_\omega\,\beta\,e_{US,t}^u,
\]
and similarly
\[
Q_{W,t}^u
=
(1-\omega_t^u)S_t^u+(1-\omega_t^{*,u})S_t^{*,u}
\ge
(1-\omega_t^u)S_t^u
\ge
\underline\lambda_\omega\,\beta\,e_{US,t}^u.
\]

For $z=b$, 
\[
Q_{US,t}^b
\ge
\omega_t^b S_t^b
\ge
\underline\lambda_\omega(1-\theta_t^b)A_t^b
\ge
\underline\lambda_\omega\,\beta\,e_{US,t}^b,
\]
and likewise
\[
Q_{W,t}^b
\ge
(1-\omega_t^b)S_t^b
\ge
\underline\lambda_\omega\,\beta \,e_{US,t}^b.
\]
Hence
\[
Q_{i,t}^z\ge \underline q\,e_{US,t}^z,
\qquad
\underline q\equiv \underline\lambda_\omega\,\beta.
\]

For the upper bound, \eqref{Qsum_identity_mod_zerosupply},
\Cref{ass_relative_dividend_bounds} imply
\[
Q_{i,t}^z
\le
Q_{US,t}^z+Q_{W,t}^z
=
\beta e_{US,t}^z+\frac{\beta+\chi}{1+\chi}e_{W,t}
\le
\left(
\beta
+\frac{\beta+\chi}{1+\chi}\overline\lambda_{e_W}
\right)e_{US,t}^z.
\]
Thus
\[
Q_{i,t}^z\le \overline q\,e_{US,t}^z,
\qquad
\overline q\equiv
\beta
+\frac{\beta+\chi}{1+\chi}\overline\lambda_{e_W}.
\]
\end{proof}

\begin{lem}[Equilibrium bounds on risky returns]\label{lem_risky_return_bounds}
Suppose
\Cref{ass_relative_dividend_bounds,ass_switch_us_fundamentals,ass_interior_equity_weights,ass_u_dividend_growth}
hold, and suppose dividends are nonnegative:
\[
D_{US,t}\ge 0,
\qquad
D_{W,t}\ge 0,
\qquad \forall t.
\]
Then there exist constants $0< \overline R<\infty$ such that, for each
$z\in\{u,b\}$ and all $t\le\tau$,
\[
0
<
R_{US,t}^z,\ R_{W,t}^z,\ R_{p,t}^z,\ R_{p,t}^{*,z}
\le
\overline R^z.
\]
\end{lem}

\begin{proof}
By \Cref{lem_stock_value_bounds},
\[
\underline q\,e_{US,t}^z
\le
Q_{i,t}^z
\le
\overline q\,e_{US,t}^z,
\qquad
i\in\{US,W\},\ z\in\{u,b\}.
\]
Let
\[
\overline d_{US}\equiv \overline\lambda_{D_{US}},
\qquad
\overline d_W\equiv \overline\lambda_{D_W},
\qquad
\overline d\equiv \max\{\overline d_{US},\overline d_W\}.
\]
Then
\[
D_{US,t}^z\le \overline d_{US}\,e_{US,t}^z,
\qquad
D_{W,t}^z\le \overline d_W\,e_{US,t}^z,
\qquad
z\in\{u,b\},\ \forall t\le\tau.
\]

For $i\in\{US,W\}$ and $z\in\{u,b\}$, the one-period stock return from date $t-1$ on the $u$-history to
date $t$ on branch $z$ is
\[
R_{i,t}^z=\frac{Q_{i,t}^z+D_{i,t}^z}{Q_{i,t-1}^u}.
\]
Since dividends are nonnegative,
\begin{align*}
\begin{split}
R_{i,t}^u
& \ge
\frac{Q_{i,t}^u}{Q_{i,t-1}^u}
\ge
\frac{\underline q\,e_{US,t}^u}{\overline q\,e_{US,t-1}^u}
\ge
\frac{\underline q\,\underline g_u}{\overline q}
\equiv
\underline R^u
>0.
\\
R_{i,t}^b
& \ge
\frac{Q_{i,t}^b}{Q_{i,t-1}^u}
\ge
\frac{\underline q\,e_{US,t}^b}{\overline q\,e_{US,t-1}^u}
\equiv
\underline R^b_t
>0.
\end{split}
\end{align*}
Similarly,
\begin{align*}
\begin{split}
R_{i,t}^u
\le
\frac{Q_{i,t}^u+D_{i,t}^u}{Q_{i,t-1}^u}
\le
\frac{(\overline q+\overline d)e_{US,t}^u}{\underline q\,e_{US,t-1}^u}
\le
\frac{(\overline q+\overline d)\overline g_u }{\underline q}
\equiv
\overline R^u
<\infty.
\\
R_{i,t}^b
\le
\frac{Q_{i,t}^b+D_{i,t}^b}{Q_{i,t-1}^u}
\le
\frac{(\overline q+\overline d)e_{US,t}^b}{\underline q\,e_{US,t-1}^u}
\le
\frac{(\overline q+\overline d)\overline g_u \overline \lambda_e}{\underline q}
\equiv
\overline R^b
<\infty.
\end{split}
\end{align*}

Now
\[
R_{p,t}^z
=
\omega_{t-1}^u R_{US,t}^z+(1-\omega_{t-1}^u)R_{W,t}^z,
\]
and
\[
R_{p,t}^{*,z}
=
\omega_{t-1}^{*,u} R_{US,t}^z+(1-\omega_{t-1}^{*,u})R_{W,t}^z.
\]
Because \Cref{ass_interior_equity_weights} implies
\[
0\le \omega_{t-1}^u,\ \omega_{t-1}^{*,u}\le 1,
\]
both $R_{p,t}^z$ and $R_{p,t}^{*,z}$ are convex combinations of
$R_{US,t}^z$ and $R_{W,t}^z$. Therefore
\[
0
<
R_{p,t}^z,\ R_{p,t}^{*,z}
\le
\overline R^z.
\]
\end{proof}

\begin{lem}[Equilibrium upper bound on the US risk-free rate]\label{lem_rf_upper_bound}
Under the assumptions of \Cref{lem_risky_return_bounds}, there exists
$\overline R_f<\infty$ such that
\[
R_{f,t}\le \overline R_f
\qquad\text{for all }t\le\tau.
\]
One may take
\[
\overline R_f=\overline R,
\]
where $\overline R\equiv \max \{\overline R^u, \overline R^b\}$ is given by \Cref{lem_risky_return_bounds}.
\end{lem}

\begin{proof}
From \eqref{RoW_thetaW_FOC_mod},
\[
\beta E_t\!\left[\mathcal M_{t,t+1}^*\bigl(R_{f,t}^W-R_{p,t+1}^*\bigr)\right]=0,
\]
so
\[
R_{f,t}^W
=
\frac{E_t\!\left[\mathcal M_{t,t+1}^*R_{p,t+1}^*\right]}
{E_t\!\left[\mathcal M_{t,t+1}^*\right]}.
\]
Since $\mathcal M_{t,t+1}^*>0$ and, by \Cref{lem_risky_return_bounds} we can let $\overline R\equiv \max \{\overline R^u, \overline R^b\}$,
\[
R_{p,t+1}^*\le \overline R
\qquad\text{a.s.},
\]
it follows that
\[
R_{f,t}^W\le \overline R.
\]
Recall that $q_{US,t}>q_{W,t}$, hence
\[
R_{f,t}=\frac{1}{q_{US,t}}<\frac{1}{q_{W,t}}=R_{f,t}^W\le \overline R.
\]
Therefore
\[
R_{f,t}\le \overline R_f,
\qquad
\overline R_f\equiv \overline R.
\]
\end{proof}


As shown in \eqref{theta_nonpos}, the convenience-yield term forces the US to be a net borrower:
\[
\theta_t<0,
\qquad
1-\theta_t>1.
\]
Thus, for the US, only the left tail $\theta_t\downarrow-\infty$ is potentially problematic. The next lemma
shows that the issuance cost prevents this from happening.

\begin{lem}[Endogenous lower bound on the US bond share]\label{lem2}
Suppose the assumptions of \Cref{lem_risky_return_bounds,lem_rf_upper_bound} hold, and suppose
\Cref{ass_issuance_costs,ass_gamma_lt_one} hold. Then there exists a constant $\underline\theta>-\infty$ such that
\[
\theta_t^u\ge \underline\theta
\qquad\text{for all }t<\tau.
\]
\end{lem}

\begin{proof}
Fix $t<\tau$, so the current state is $u$. After saving separation, the US chooses $(\omega_t,\theta_t)$ to maximize
\[
J_t(\omega,\theta)
\equiv
\frac{\beta}{1-\gamma}\log E_t\!\left[R_{A,t+1}(\omega,\theta)^{1-\gamma}\right]
-\frac{\kappa}{2}(\omega-\bar\omega)^2
-\eta\,\Upsilon(\theta),
\]
where
\[
R_{A,t+1}(\omega,\theta)
=
(1-\theta)R_{p,t+1}(\omega)+\theta R_{f,t},
\qquad
R_{p,t+1}(\omega)=\omega R_{US,t+1}+(1-\omega)R_{W,t+1}.
\]

Consider the feasible reference choice $(\omega,\theta)=(\bar\omega,0)$. By \Cref{lem_risky_return_bounds},
\[
R^u_{p,t+1}(\bar\omega) \ge \underline R^u >0.
\]
Because $0<\gamma<1$ and the $u$ branch has probability $\pi$, this choice
yields
\[
E_t\!\left[R_{p,t+1}(\bar\omega)^{1-\gamma}\right]
\ge
\pi(\underline R^u)^{1-\gamma},
\]
and hence a finite lower bound
\[
J_t(\bar\omega,0)\ge \underline J>-\infty,
\]
where $\underline J$ is independent of $t$.

Now evaluate the objective at the equilibrium choice $(\omega_t,\theta_t)$. Since $\theta_t<0$ by
\eqref{theta_nonpos}, and since \Cref{lem_risky_return_bounds,lem_rf_upper_bound} give
\[
R_{p,t+1}(\omega_t)\le \overline R,
\qquad
R_{f,t}\le \overline R_f,
\]
we have
\[
R_{A,t+1}(\omega_t,\theta_t)
=
(1-\theta_t)R_{p,t+1}(\omega_t)+\theta_t R_{f,t}
\le
(1-\theta_t)\overline R
\le
\overline R(1+|\theta_t|).
\]
Hence
\[
\frac{\beta}{1-\gamma}\log E_t\!\left[R_{A,t+1}(\omega_t,\theta_t)^{1-\gamma}\right]
\le
\beta\log\!\bigl(\overline R(1+|\theta_t|)\bigr),
\]
and therefore
\[
J_t(\omega_t,\theta_t)
\le
\beta\log\!\bigl(\overline R(1+|\theta_t|)\bigr)-\eta\,\Upsilon(\theta_t).
\]
Since $(\omega_t,\theta_t)$ is optimal,
\[
J_t(\omega_t,\theta_t)\ge J_t(\bar\omega,0)\ge \underline J.
\]
Thus
\[
\beta\log\!\bigl(\overline R(1+|\theta_t|)\bigr)-\eta\,\Upsilon(\theta_t)\ge \underline J.
\]

By \Cref{ass_issuance_costs},
\[
\lim_{\theta\downarrow -\infty}\Upsilon'(\theta)=-\infty,
\]
so, in particular,
\[
\Upsilon(\theta)\to+\infty
\qquad\text{as }\theta\downarrow -\infty,
\]
and $\Upsilon(\theta)$ dominates $\log(1+|\theta|)$ in the left tail. Hence the last inequality cannot hold if
$\theta_t\downarrow -\infty$. Therefore there exists $\underline\theta>-\infty$ such that
\[
\theta_t^u\ge \underline\theta
\qquad\text{for all }t<\tau.
\]
\end{proof}


\begin{ass}[Risk-free return bounded away from zero]\label{ass_rf_lower_bound}
There exists $\underline R_f>0$ such that
$$
R_{f,t}\ge \underline R_f
\qquad\text{for all }t\le \tau.
$$
\end{ass}

\begin{lem}[Uniform slack from the bad-state natural debt limit]
\label{lem_uniform_natural_debt_slack}
Suppose the assumptions of
\Cref{lem_risky_return_bounds,lem_rf_upper_bound,lem2}
hold. Suppose also that \Cref{ass_rf_lower_bound} holds and that
$R^u_{A,t}$ is bounded above and bounded away from zero. Suppose the equilibrium
choice is interior in $\theta_t$, so that the bond-share FOC holds with equality.
Then there exists
$\underline\delta\in(0,1)$ such that, whenever $z_{t-1}=u$,
$$
R^b_{A,t}\ge \underline\delta R^b_{p,t}.
$$
\end{lem}
\begin{proof}
Let
\[
r_t\equiv R^b_{p,t},\qquad
x_t\equiv R_{f,t-1}-R^b_{p,t},
\qquad
\delta_t\equiv \frac{R^b_{A,t}}{R^b_{p,t}}.
\]
Since
\[
R^b_{A,t}
=
r_t+\theta^u_{t-1}x_t,
\]
if $x_t\le0$, then $\theta^u_{t-1}<0$ implies
\[
R^b_{A,t}\ge r_t,
\]
and hence $\delta_t\ge1$. Therefore only the case $x_t>0$ is relevant.

Because admissible portfolio choices must deliver strictly positive old-age
consumption in every successor state, and because \(A_{t-1}^u>0\), we have
\(R^b_{A,t}>0\). Hence \(\delta_t>0\). Suppose, toward contradiction, that along some subsequence $\delta_t\to0$.
Then by
\[
\delta_t r_t
=
r_t+\theta^u_{t-1}x_t,
\]
we have that
\[
\theta^u_{t-1}
=
-\frac{(1-\delta_t)r_t}{x_t}.
\]
Since $\delta_t\to0$, for all sufficiently large $t$,
\[
1-\delta_t\ge \frac12.
\]
Let $B\equiv-\underline\theta<\infty$. Because
$\theta^u_{t-1}\in[\underline\theta,0)$,
\[
\frac{(1-\delta_t)r_t}{x_t}
=
|\theta^u_{t-1}|
\le B.
\]
Therefore,
\[
r_t\le 2B x_t.
\]
Since
\[
R_{f,t-1}=r_t+x_t
\]
and $R_{f,t-1}\ge \underline R_f>0$, it follows that
\[
x_t
\ge
\frac{\underline R_f}{2B+1}
>0.
\]
Moreover, because $x_t>0$, $r_t<R_{f,t-1}\le \overline R_f$ by
\Cref{lem_rf_upper_bound}, so
\[
0\le R^b_{A,t}
=
\delta_t r_t
\le
\delta_t \overline R_f
\to0.
\]

The bond-share FOC is
\[
\beta
E_{t-1}\left[
\mathcal K_{t-1,t}(R_{f,t-1}-R_{p,t})
\right]
=
\eta\Upsilon'(\theta^u_{t-1}).
\]
The right-hand side is uniformly bounded because
$\theta^u_{t-1}\in[\underline\theta,0]$ and $\Upsilon'$ is continuous.
The $u$-branch contribution is uniformly bounded because $R^u_{A,t}$ is
bounded above and bounded away from zero, while $R_{f,t-1}$ and $R^u_{p,t}$
are uniformly bounded by \Cref{lem_risky_return_bounds,lem_rf_upper_bound}.

On the $b$ branch,
\[
\mathcal K^b_{t-1,t}
=
\frac{(R^b_{A,t})^{-\gamma}}
{\pi (R^u_{A,t})^{1-\gamma}
+(1-\pi)(R^b_{A,t})^{1-\gamma}}.
\]
Since $R^b_{A,t}\to0$ and $R^u_{A,t}$ is bounded above and bounded away from
zero, we have
\[
\mathcal K^b_{t-1,t}\to\infty.
\]
Together with
\[
x_t\ge \frac{\underline R_f}{2B+1}>0,
\]
this implies
\[
\mathcal K^b_{t-1,t}x_t\to\infty.
\]
Thus the $b$-branch contribution to the left-hand side of the FOC diverges to
$+\infty$, while the $u$-branch contribution and the right-hand side are uniformly
bounded. This contradiction implies that $\delta_t$ is bounded away from zero.
Hence there exists $\underline\delta\in(0,1)$ such that
\[
R^b_{A,t}\ge \underline\delta R^b_{p,t}.
\]
\end{proof}

\begin{ass}[Positive lower bound on continuation total-saving returns]\label{ass_positive_u_total_saving_return}
\[
(1-\underline\theta)\underline R^u+\underline\theta\,\overline R_f>0.
\]
where $\underline\theta$ is given by \Cref{lem2}, and
$\underline R^u,\overline R_f$ are given by
\Cref{lem_risky_return_bounds,lem_rf_upper_bound}.
\end{ass}

\begin{lem}[Equilibrium bounds on total-saving returns and old-age consumption]\label{lem_total_return_consumption_bounds}
Suppose
\Cref{ass_issuance_costs,ass_gamma_lt_one,ass_rf_lower_bound,ass_relative_dividend_bounds,ass_switch_us_fundamentals,ass_interior_equity_weights,ass_u_dividend_growth,ass_positive_u_total_saving_return}
hold.
Then there exist constants $0< \overline R^z_A<\infty$ and
$0<\underline\lambda^z_C\le \overline\lambda^z_C<\infty$ such that, for each $z\in\{u,b\}$ and all $t\le\tau$,
\[
0< R_{A,t}^z\le \overline R^z_A<\infty,
\]
and
\[
\underline\lambda^z_C\,e_{US,t}^z
\le
C_t^z
\le
\overline\lambda^z_C\,e_{US,t}^z.
\]
\end{lem}

\begin{proof}
\Cref{lem_risky_return_bounds,lem_rf_upper_bound,lem2} imply that there exist constants
$0<\underline R^u\le \overline R^u<\infty$, $\overline R^b<\infty$,
$\overline R_f<\infty$, and $\underline\theta>-\infty$, together with
$\underline R^b_t>0$, such that
\[
\underline R^u\le R_{p,t}^u\le \overline R^u,
\quad
0<\underline R^b_t\le R_{p,t}^b\le \overline R^b,
\quad
R_{f,t}\le \overline R_f,
\quad
	\theta_{t-1}^u\ge \underline{\theta},
\quad
z\in\{u,b\},\ \forall t\le\tau.
\]
Also, by \eqref{theta_nonpos},
\[
	\theta_{t-1}^u<0
\qquad\text{for all }t\le\tau.
\]

For $z\in\{u,b\}$,
\[
R_{A,t}^z=(1-\theta_{t-1}^u)R_{p,t}^z+\theta_{t-1}^u R_{f,t-1}.
\]
Using
\[
R_{p,t}^u\ge \underline R^u,
\qquad
R_{f,t-1}\le \overline R_f,
\qquad
	\theta_{t-1}^u\in[\underline\theta,0),
\]
we obtain
\[
R_{A,t}^u
\ge
(1-\theta_{t-1}^u)\underline R^u+\theta_{t-1}^u\overline R_f
\ge
\min\!\left\{
\underline R^u,\,
(1-\underline\theta)\underline R^u+\underline\theta\,\overline R_f
\right\}
\equiv
\underline R^u_A.
\]
By \Cref{ass_positive_u_total_saving_return}, $\underline R^u_A>0$.

Similarly,
\[
R_{A,t}^z
\le
(1-\theta_{t-1}^u)\overline R^z+\theta_{t-1}^u R_{f,t-1}
\le
(1-\theta_{t-1}^u)\overline R^z
\le
(1-\underline\theta)\overline R^z
\equiv
\overline R^z_A.
\]
Therefore,
\begin{align*}
    0<\underline R^u_{A}\le R_{A,t}^u\le \overline R^u_A<\infty,
    \\
    0<R_{A,t}^b\le \overline R^b_A<\infty.
\end{align*}
The strict positivity of $R^b_{A,t}$ follows from admissibility of old-age
consumption in every successor state and $A_{t-1}^u>0$.

% Now show C^u is proportional to e^u
Since
\[
C_t^z=A_{t-1}^uR_{A,t}^z=\beta e_{US,t-1}^uR_{A,t}^z,
\]
and by \Cref{ass_u_dividend_growth} it follows that
\[
\frac{\beta\underline R^u_A}{\overline g_u}\,e_{US,t}^u
\le 
\beta e_{US,t-1}^u\underline R^u_A
\le
C_t^u
\le
\beta e_{US,t-1}^u\overline R^u_A
\le
\frac{\beta\overline R^u_A}{\underline g_u}\,e_{US,t}^u.
\]
% Show C^b is proportional to e^b
For $C^b_t$, by \Cref{lem_uniform_natural_debt_slack},
\[
R^b_{A,t}\ge \underline\delta R^b_{p,t}.
\]
Therefore,
\[
C_t^b
=
\beta e^u_{US,t-1}R^b_{A,t}
\ge
\beta \underline\delta e^u_{US,t-1}R^b_{p,t}.
\]
Using the risky-return lower bound,
\[
R^b_{p,t}
\ge
\frac{\underline q}{\overline q}
\frac{e^b_{US,t}}{e^u_{US,t-1}},
\]
we obtain
\[
C_t^b
\ge
\beta\underline\delta
\frac{\underline q}{\overline q}
e^b_{US,t}.
\]
For the upper bound, since $\theta^u_{t-1}<0$ and $R_{f,t-1}>0$,
\[
R^b_{A,t}
=
(1-\theta^u_{t-1})R^b_{p,t}
+\theta^u_{t-1}R_{f,t-1}
\le
(1-\theta^u_{t-1})R^b_{p,t}
\le
(1-\underline\theta)R^b_{p,t}.
\]
Using the risky-return upper bound,
\[
R^b_{p,t}
\le
\frac{(\overline q+\overline d)e^b_{US,t}}
{\underline q e^u_{US,t-1}},
\]
we obtain
\[
C_t^b
=
\beta e^u_{US,t-1}R^b_{A,t}
\le
\beta(1-\underline\theta)
\frac{\overline q+\overline d}{\underline q}
e^b_{US,t}.
\]

Hence one may take
\begin{align*}
    &\underline\lambda^u_C\equiv \frac{\beta\underline R^u_A}{\overline g_u}>0,
\qquad
\overline\lambda^u_C\equiv \frac{\beta\overline R^u_A}{\underline g_u}<\infty,
\\ &
\underline\lambda^b_C
\equiv
\beta\underline\delta\frac{\underline q}{\overline q}>0.
\\ &
\overline\lambda^b_C
\equiv
\beta(1-\underline{\theta})\frac{\overline q+\overline d}{\underline{q}}<\infty,
\end{align*}
which proves the claim.
\end{proof}


With these ingredients, we can state a single
corollary that gives primitive conditions ensuring that the summability conditions
in \Cref{thm_1} reduce to \eqref{thm_cond2}.

\begin{coro}\label{coro1}
Suppose
\Cref{ass_regular_kernel,ass_issuance_costs,ass_gamma_lt_one,ass_rf_lower_bound,ass_regime_switch,ass_balanced_growth,ass_switch_dependence,ass_relative_dividend_bounds,ass_switch_us_fundamentals,ass_interior_equity_weights,ass_u_dividend_growth,ass_positive_u_total_saving_return}
hold. 
For $z\in\{u,b\}$, let $x_t^z$ be the value of $x_{t}^z$ conditional on $z_0=\cdots=z_{t-1}=u$ and $z_t=z$.
For notational simplicity, write \(e_t^z=e_{US,t}^z\),
\(D_t^z=D_{US,t}^z\), and \(Q_t^z=Q_{US,t}^z\).
Then the two conditions \eqref{thm_cond1} in \Cref{thm_1} are equivalent to the primitive conditions
\begin{subequations}\label{thm_cond2}
\begin{align}
\sum_{t=1}^\infty \frac{D_t^u}{e_{t}^u}&<\infty, \label{thm_cond_2a}\\
\sum_{t=1}^\infty
\left(\frac{e_{t}^b}{e_{t}^u}\right)^{1-\gamma}&<\infty. \label{thm_cond_2b}
\end{align}
\end{subequations}
\end{coro}

\begin{proof}
We show (a) \eqref{thm_cond_1a} $\Longleftrightarrow$ \eqref{thm_cond_2a} and
(b) \eqref{thm_cond_1b} $\Longleftrightarrow$ \eqref{thm_cond_2b}. 

\medskip
\noindent\textbf{(a) Proof of \eqref{thm_cond_1a} $\Longleftrightarrow$ \eqref{thm_cond_2a}.}
By \Cref{lem_stock_value_bounds}, there exist constants $0<\underline q\le \overline q<\infty$ such that
\[
\underline q\,e_{t}^u\le Q_{t}^u\le \overline q\,e_{t}^u
\qquad\text{for all }t\le\tau.
\]
Therefore
\[
\sum_{t=1}^\infty \frac{D_t^u}{Q_t^u}<\infty
\quad\Longleftrightarrow\quad
\sum_{t=1}^\infty \frac{D_t^u}{e_{t}^u}<\infty.
\]

\medskip
\noindent\textbf{(b) Proof of \eqref{thm_cond_1b} $\Longleftrightarrow$ \eqref{thm_cond_2b}.}

By \Cref{lem_stock_value_bounds},
\[
\underline q\,e_{t}^z\le Q_t^z\le \overline q\,e_{t}^z
\qquad\text{for all }t\le\tau.
\]
By \Cref{ass_u_dividend_growth}, there exist constants $0<\underline q_b\le \overline q_b<\infty$ such that
\[
\underline q_b\,e_{t}^b
\le
Q_t^b+D_t^b
\le
\overline q_b\,e_{t}^b
\qquad\text{for all }t\le\tau,
\]
where one may take
\[
\underline q_b\equiv \underline q,
\qquad
\overline q_b\equiv \overline q+\overline\lambda_{D_{US}}.
\]

By
\Cref{lem_total_return_consumption_bounds}, there exist constants
$0<\underline\lambda^z_C\le \overline\lambda^z_C<\infty$ such that
\[
\underline\lambda^z_C\,e_{t}^z
\le
C_t^z
\le
\overline\lambda^z_C\,e_{t}^z,
\qquad
z\in\{u,b\},\ \forall t\le\tau.
\]
Hence the preceding bounds imply the two-sided comparison
\[
\frac{
\frac{Q^b_t+D^b_t}{(C^b_t)^\gamma}
}{
\frac{Q^u_t}{(C^u_t)^\gamma}
}
\asymp
\frac{e^b_t/(e^b_t)^\gamma}{e^u_t/(e^u_t)^\gamma}
=
\left(\frac{e^b_t}{e^u_t}\right)^{1-\gamma}.
\]
Therefore,
\[
\sum_{t=1}^{\infty}
\frac{
\frac{Q^b_t+D^b_t}{(C^b_t)^\gamma}
}{
\frac{Q^u_t}{(C^u_t)^\gamma}
}
<\infty
\quad\Longleftrightarrow\quad
\sum_{t=1}^{\infty}
\left(\frac{e^b_t}{e^u_t}\right)^{1-\gamma}
<\infty.
\]

This proves \eqref{thm_cond_1b} $\Longleftrightarrow$ \eqref{thm_cond_2b}.
\end{proof}

%========================================================
% Appendix
%========================================================


\appendix
%\section{Proof of Lemma~\ref{lem1}}


\bibliography{references.bib} 

\end{document}
